\section{Example: Real Beams, Part II} \label{ex9}

In example \ref{ex7}, we've made the aquaintance of a delicate set-up.
Due to the huge number of undulator periods, we are very sensitive to the
input parameters of the beam. Therefore, the assumption, that the
spectral and spatial distributions of the radiation is more or less the
same for all electrons, is not valid anymore. Hence the folding
procedures to take the energy-spread and emittance into account would
fail. To overcome this difficulties, we must track many electrons
according to the beam phase space distribution.

We consider a long beamline with $60$ undulators.
To calculate the spectrum of the spontaneous synchrotron radiation, we enter

\cboxs{\senv \grau \\
python3 ../python/wave\_example\_9.py\\
cp \$WAVE/input/wave\_example\_9.in wave.in\\
cp \$WAVE/input/wave\_example\_9\_magseq.in magseq.in \\
cp \$WAVE/input/wave\_example\_9\_wave.bmask wave.bmask\\
wave\\
wplot \\
}

This run takes some minutes, since for the considered $180\,m$ long undulator
section we need about a million integration steps for the spectrum calculation.
For this single electron beam, we get a sharp peak for the on-axis
flux-density (Fig. \ref{ex9fig1}).

\begin{figure}[htb]
    \centering
    %\framebox(12,9){
      \includegraphics[width=1\textwidth]
      {/home/scheer/wav/doc/pp/wave_example_9_fig1.pdf}
    %} % endframebox
    \caption{
      5th harmonic of a single electron beam for 30
      undulator sections including continous energy-loss and tapering of
      the undulators
    }
    \label{ex9fig1}
\end{figure}

The set-up of one six meter subsection section including two
undulators is given in the file \scap{magseq.in}. In combination with\\

\cboxss{XSTART=0.0\\
XSTOP=180.\\
IPERIODG=1\\
PERIODG=6.\\}

this results in a beamline of $180\,m$ length, where the field strength is
corrected for the energy loss due to the spontaneous synchrotron radiation.
Each section includes phase-shifters and quadrupoles to match the
phase-advance and $\beta$-functions of the $30$ sections.

Note, that for time reasons, only a quarter of flux is calculated,
i.e. the pinhole is shifted, such that the lower left corner is
on-axis. This is achieved by \\

\cboxss{IPBRILL=1\\
PINCEN(2)=9999.\\
PINCEN(3)=9999.\\}

Now, we come to the actual task to calculate the $5th$ harmonic for a
real beam, including energy-spread, emittance, and an energy-loss with
quantum fluctuations. We start step by step setting in \wi \\

\cboxss{IENELOSS=-1\\
IBUNCH=-1\\
NBUNCH=1000\\
NEINBUNCH=1\\
IPIN=0\\
}

The variable \scap{IBUNCH} triggers the bunch mode, i.e. to track a
given number of bunches and electrons per bunch (\scap{NBUNCH} and
\scap{NEINBUNCH} respectively). The radiation of different bunches is
added incoherently, while the contributions of electrons of the same
bunch are treated coherently. Thus, the setting given above calculates
the radiation of $1000$ electron, summed up incoherently. The
resulting spectra are normalized to the beam current. Please, refer to
the namelist \scap{BUNCHN} for further options.

We run \W and get (after some minutes) the results shown in Fig. \ref{ex9fig2}. We see some
phase-space distributions of the beam electrons and the corresponding
flux-density spectrum of a single electron beam in the lower right
corner together with the same spectrum for the real beam.

We clearly see the flux-density reduced by the emittance and
energy-spread. It reaches only about $0.5 \%$ of the original
one.\\[3cm]

\begin{figure}[htb]
    \centering
    %\framebox(12,9){
      \includegraphics[width=0.9\textwidth]
      {/home/scheer/wav/doc/pp/wave_example_9_fig2.pdf}
    %} % endframebox
    \caption{
      Flux-density of the 5th harmonic of $1000$ electrons beam for 30
      undulator sections including energy-loss with quantum
      fluctuations and tapering of
      the undulators (lower right plot)
    }
    \label{ex9fig2}
\end{figure}

\vspace*{1cm}
\underline{Flux through a pinhole}\\

Since the calculation of radiation distributions in a pinhole are very
time consuming, the mode \scap{IPIN=3} can be used. In this mode, the
radiation is not calculated for all grid points of the pinhole, but
only for one randomly chosen point for each electron track. The idea
is, that the random phase-space coordinates of the electron and the
random observation point are not correlated, thus, for large enough a
number of tracks, the flux trough the pinhole can be calculated by
sampling the phase-space and the observation point simultaneously.

Edit \wi~and set:

\cboxss{IPIN=3\\
IPBRILL=0
PINW=0.00014\\
PINH=0.00014}

and run \W and get Fig. \ref{ex9fig3}.

\begin{figure}[htb]
    \centering
    %\framebox(12,9){
      \includegraphics[width=0.9\textwidth]
      {/home/scheer/wav/doc/pp/wave_example_9_fig3.pdf}
    %} % endframebox
    \caption{
      Flux of the 5th harmonic of $1000$ electrons beam for 30
      undulator sections including energy-loss with quantum
      fluctuations and tapering of
      the undulators (lower right plot)
    }
    \label{ex9fig3}
\end{figure}

The flux is also significantly reduced compared to that of zero
emittance and energy spread beam (Fig. \ref{ex9fig1}).

%\newpage
\underline{Proposed exercises:}

\begin{itemize}

\item
Copy \we to \wie, run \W, and plot the results. Store the results and try
to run \W in the multi-particle mode to compare with the previous
results.

\item
Check the namelist \scap{\$BUNCHN} in \wi and try to provide
your own phasespace distribution.

\item
Check the namelist \scap{\$BUNCHN} in \wie, and try to
calculate coherent radiation of electrons.

\end{itemize}

